% ============================================================
% 学习系统：如何快速进入一个新领域
% ============================================================

\section{领域学习操作系统}

\begin{conceptbox}
\textbf{目标：} 学一个新领域时，不只是收集资料，而是快速建立三层能力：
\begin{itemize}[leftmargin=1.5em]
  \item \key{地图感}：知道这个领域有哪些问题、方法、流派和关键论文。
  \item \key{原理感}：能把核心方法还原成目标函数、数据流、训练流程和工程约束。
  \item \key{判断力}：能判断一个新方法解决了什么痛点、代价是什么、是否真的重要。
\end{itemize}
\end{conceptbox}

\subsection{核心循环}

学习一个领域时，固定走下面六步：

\begin{enumerate}[leftmargin=1.5em]
  \item \textbf{建地图：} 先列出 5--10 个核心问题，而不是先读论文。
  \item \textbf{抓主干：} 找到最小闭环：输入是什么，输出是什么，训练信号是什么，优化什么。
  \item \textbf{读经典：} 每条路线只选 1--2 篇代表作，先读懂方法结构。
  \item \textbf{拆公式：} 把损失函数、采样流程、奖励定义拆成自己的符号系统。
  \item \textbf{做复现：} 不一定完整复现，至少复现一个 toy example、loss 计算或流程图。
  \item \textbf{写判断：} 最后用自己的话写：它解决了什么、失败在哪里、下一步会怎么演化。
\end{enumerate}

\begin{summarybox}
\textbf{一句话：} 领域学习不是从资料开始，而是从\key{问题结构}开始；不是以“看完多少论文”为目标，而是以“能解释、能比较、能预测”为目标。
\end{summarybox}

\subsection{三遍阅读法}

\textbf{第一遍：地图阅读。}
只看标题、摘要、图 1、方法总览和实验表，目标是判断它属于哪条技术路线。

\textbf{第二遍：机制阅读。}
重点看方法章节，提取：
\begin{itemize}[leftmargin=1.5em]
  \item 优化目标是什么；
  \item 数据从哪里来；
  \item 训练时哪些模型参与；
  \item 推理时和原方法相比多了什么。
\end{itemize}

\textbf{第三遍：批判阅读。}
问四个问题：
\begin{itemize}[leftmargin=1.5em]
  \item 这个方法真正的新东西是什么？
  \item 它绕开了什么困难，又引入了什么代价？
  \item 实验是否能支撑论文的主张？
  \item 如果迁移到我的方向，会卡在哪里？
\end{itemize}

\subsection{每天的学习产物}

每天至少留下四种产物之一：

\begin{itemize}[leftmargin=1.5em]
  \item \textbf{一张地图：} 方法谱系、模块关系、训练流程。
  \item \textbf{一张表：} 方法对比、假设对比、损失函数对比。
  \item \textbf{一个公式拆解：} 从原论文公式翻译成自己的变量解释。
  \item \textbf{一个最小实验：} toy code、伪代码、loss 计算、数据流模拟。
\end{itemize}

\begin{warnbox}
\textbf{避免低效学习：} 不要把“读了很多论文”等价于“懂了这个领域”。真正的掌握标准是：给你一篇新论文，你能很快判断它属于哪条路线、解决了哪个旧问题、核心假设是否可靠。
\end{warnbox}
